\section{The ZISAR Model}\label{sec:model}

% ---------------------------------------------------------------------------
\subsection{Notation and setup}\label{sec:model_notation}

Consider $n$ spatial units indexed by $i = 1,\ldots,n$, each associated with a non-negative integer outcome $y_i \in \{0,1,2,\ldots\}$.
Let $\bW = (w_{ij})_{n \times n}$ be a known, row-standardised spatial weights matrix with $w_{ii}=0$ for all $i$; typical choices include queen or rook contiguity on a regular lattice, or $k$-nearest-neighbour graphs on irregular point patterns \citep{anselin1988,arbia2006}.
Each unit is endowed with a $(p \times 1)$ vector of exogenous covariates $\mathbf{x}_i$ entering the count process and a $(q \times 1)$ vector $\mathbf{z}_i$ entering the zero-inflation process (with $\mathbf{z}_i$ possibly overlapping with or equal to $\mathbf{x}_i$).
Stacking observations, we write $\by = (y_1,\ldots,y_n)'$, $\bX = (\mathbf{x}_1,\ldots,\mathbf{x}_n)'$, and $\bZ = (\mathbf{z}_1,\ldots,\mathbf{z}_n)'$.
Let $\lambda_1 \geq \ldots \geq \lambda_n$ denote the eigenvalues of $\bW$.

% ---------------------------------------------------------------------------
\subsection{Model specification}\label{sec:model_spec}

The Zero-Inflated Spatial Autoregressive (ZISAR) model posits that each observation $y_i$ is drawn from a two-component mixture:
\begin{equation}\label{eq:mixture}
  y_i \sim
  \begin{cases}
    0 & \text{with probability } \pi_i, \\
    \mathrm{NB}(\mu_i, \kappa) & \text{with probability } 1-\pi_i,
  \end{cases}
\end{equation}
where $\mathrm{NB}(\mu_i,\kappa)$ denotes the negative binomial distribution with conditional mean $\mu_i > 0$ and shape parameter $\kappa > 0$, so that $\Var[y_i \mid s_i=0] = \mu_i + \mu_i^2/\kappa$.
The probability $\pi_i \in (0,1)$ governs the zero-inflation (or ``structural zero'') mechanism.

The defining feature of ZISAR---what distinguishes it from standard zero-inflated models and from CAR-based spatial alternatives---is that the conditional mean of the count process is governed by a \emph{latent spatial autoregressive field}.

\paragraph{Count process: SAR on the latent predictor.}
Define the latent linear predictor $\boldsymbol{\eta} = (\eta_1,\ldots,\eta_n)'$ through the simultaneous autoregressive system
\begin{equation}\label{eq:sar_eta}
  \boldsymbol{\eta} = \rho\,\bW\boldsymbol{\eta} + \bX\bbeta + \bu, \qquad \bu \sim N(\mathbf{0},\,\sigma^2\mathbf{I}_n),
\end{equation}
which yields the reduced form
\begin{equation}\label{eq:reduced_form}
  \boldsymbol{\eta} = (\mathbf{I}_n - \rho\,\bW)^{-1}(\bX\bbeta + \bu) = \mathbf{S}(\rho)\,(\bX\bbeta + \bu),
\end{equation}
where $\mathbf{S}(\rho) \equiv (\mathbf{I}_n - \rho\,\bW)^{-1}$ is the \emph{spatial multiplier matrix}.
The conditional mean of the count process is then
\begin{equation}\label{eq:count_mean}
  \mu_i = \exp(\eta_i), \qquad i = 1,\ldots,n.
\end{equation}

The scalar parameter $\rho$ governs the strength and range of spatial interaction: a change in $\bX\bbeta$ in unit $j$ propagates to $\eta_i$ for all $i$ through the power-series expansion
\begin{equation}\label{eq:power_series}
  \mathbf{S}(\rho) = \sum_{s=0}^{\infty} \rho^s\,\bW^s,
\end{equation}
which converges for $|\rho| < 1/\max(|\lambda_1|,|\lambda_n|)$.
The random effect $u_i$ captures unit-specific unobserved heterogeneity in the count intensity, and is absorbed by the spatial multiplier: the implied covariance of $\boldsymbol{\eta}$ is
\begin{equation}\label{eq:cov_eta}
  \Var[\boldsymbol{\eta}] = \sigma^2\,\mathbf{S}(\rho)\,\mathbf{S}(\rho)' = \sigma^2\bigl[(\mathbf{I}_n-\rho\bW)'(\mathbf{I}_n-\rho\bW)\bigr]^{-1},
\end{equation}
a simultaneous autoregressive covariance structure \citep{anselin1988} that should be contrasted with the conditional autoregressive (CAR) specification used in disease-mapping models \citep{besag_york_mollie1991,lee2013}.

\begin{remark}[SAR on $\boldsymbol{\eta}$ versus SAR on $\by$]\label{rem:sar_eta_vs_y}
An alternative specification places the spatial lag on the observed count: $\log(\mu_i) = \mathbf{x}_i'\bbeta + \rho\sum_j w_{ij}y_j + v_i$.
While this ``SAR-on-$\by$'' formulation has a direct contagion interpretation, it introduces endogeneity---$y_j$ is simultaneously an outcome and a regressor---requiring instrumental-variable workarounds \citep{lambert2010}.
The SAR-on-$\boldsymbol{\eta}$ formulation \eqref{eq:sar_eta} avoids this: the spatial multiplier operates on the \emph{latent} predictor, not on the observed count, enabling proper maximum-likelihood estimation without instruments.
Crucially, the spatial econometrics interpretation is preserved: $\mathbf{S}(\rho)$ generates the same direct/indirect/total effects decomposition as in the linear SAR \citep{lesage_pace2009}.
\end{remark}

\paragraph{Zero-inflation process.}
The structural-zero mechanism is governed by a latent Bernoulli indicator $s_i \sim \text{Bernoulli}(\pi_i)$, with
\begin{equation}\label{eq:zeroinfl_general}
  \logit(\pi_i) = \mathbf{z}_i'\bgamma + g_i.
\end{equation}
In the \emph{general ZISAR}, the spatial field $\mathbf{g} = (g_1,\ldots,g_n)'$ follows a separate SAR process:
\begin{equation}\label{eq:sar_g}
  \mathbf{g} = \rho_0\,\bW\mathbf{g} + \boldsymbol{\varepsilon}_0, \qquad \boldsymbol{\varepsilon}_0 \sim N(\mathbf{0},\,\sigma_0^2\mathbf{I}_n),
\end{equation}
yielding $\mathbf{g} = (\mathbf{I}_n - \rho_0\bW)^{-1}\boldsymbol{\varepsilon}_0$ with its own autoregressive parameter $\rho_0$.
This allows spatial clustering of the zero-inflation regime---for instance, neighbouring municipalities may all lack demand for a specific health service, and this absence is itself spatially structured.

\paragraph{Reduced-rank inflation field.}
A field with one element per unit is the natural specification, but each $g_i$ is informed by a single outcome that is zero or not, and Section~\ref{sec:est_laplace} shows that the Laplace approximation to its likelihood fails.
We therefore also consider a field of lower dimension,
\begin{equation}\label{eq:rr_g}
  \mathbf{g} = \mathbf{B}\boldsymbol{\alpha}, \qquad
  \boldsymbol{\alpha} = \rho_0\,\bW_K\boldsymbol{\alpha} + \boldsymbol{\varepsilon}_0, \qquad
  \boldsymbol{\varepsilon}_0 \sim N(\mathbf{0},\,\sigma_0^2\mathbf{I}_K),
\end{equation}
where $\boldsymbol{\alpha}$ holds $K \ll n$ values at the knots of a regular lattice laid over the study region, $\bW_K$ is the row-standardised contiguity matrix of the knots and $\mathbf{B}$ is an $n\times K$ matrix of local basis functions: entry $(i,k)$ is the bisquare function $(1 - d_{ik}^2/r^2)^2$ of the distance $d_{ik}$ between unit $i$ and knot $k$, zero beyond a radius $r$, and the rows are normalised to sum to one, so that $g_i$ is a weighted mean of the knot values around unit $i$ \citep{cressie_johannesson2008}.
Each knot value is informed by all the units within $r$ of the knot---some 150 in the simulations of Appendix~\ref{app:general}, with knots 10~km apart and $r = 15$~km---and the field is smooth at scales below the knot spacing.
The parameters are those of the general model, $\rho_0$ and $\sigma_0^2$, now describing the process at the knots, with $\rho_0 \ge 0$ imposed since a negatively autocorrelated zero regime has no interpretation.
The specification is the fixed-rank representation of spatial statistics \citep{cressie_johannesson2008} applied to the inflation predictor; its purpose here is inferential rather than computational, since the sparse algebra of Section~\ref{sec:estimation} handles the unit-level field without difficulty.

In the \emph{restricted ZISAR}, the inflation process has no spatial random effect:
\begin{equation}\label{eq:zeroinfl_restricted}
  g_i \equiv 0 \quad \forall\,i, \qquad\text{so}\qquad \logit(\pi_i) = \mathbf{z}_i'\bgamma.
\end{equation}
The restricted model is simpler, avoids identification issues between $\rho$ and $\rho_0$, and is the specification on which we focus the estimation algorithm, the Monte~Carlo study and the first empirical application (Section~\ref{sec:solar}).
The general model is examined by simulation in Appendix~\ref{app:general}: with the unit-level field the Laplace approximation fails, as Section~\ref{sec:est_laplace} explains, whereas the reduced-rank field is recovered.

% ---------------------------------------------------------------------------
\subsection{Full parameter vector and special cases}\label{sec:model_cases}

The complete parameter vector of the general ZISAR is
\begin{equation}\label{eq:theta}
  \btheta = (\bbeta', \rho, \bgamma', \rho_0, \kappa, \sigma^2, \sigma_0^2)'.
\end{equation}
For the restricted model, $\rho_0$ and $\sigma_0^2$ are dropped; with the reduced-rank field \eqref{eq:rr_g} they describe the knot process.
Table~\ref{tab:special_cases} lists the special cases nested within the specification.

\begin{table}[htbp]
\centering
\caption{Special cases nested within the ZISAR model.}\label{tab:special_cases}
\footnotesize
\begin{tabular}{lcl}
\toprule
Special case & Restrictions & Reference \\
\midrule
Restricted ZISAR & $g_i = 0$ & this paper (focus) \\
General ZISAR, reduced-rank field & $\mathbf{g} = \mathbf{B}\boldsymbol{\alpha}$, $K \ll n$ & this paper \\
SAR-NB-RE & $\pi_i = 0\;\forall\,i$ & -- \\
SAR-Poisson & $\pi_i = 0,\;\kappa\to\infty,\;\sigma^2 = 0$ & \citet{lambert2010} \\
ZINB (aspatial) & $\rho = 0,\;\sigma^2 = 0$ & \citet{cameron_trivedi2013} \\
NB-RE (aspatial) & $\pi_i = 0,\;\rho = 0$ & \citet{mcculloch_searle2001} \\
Poisson regression & $\pi_i = 0,\;\rho = 0,\;\kappa\to\infty,\;\sigma^2 = 0$ & textbook \\
\bottomrule
\end{tabular}
\end{table}

% ---------------------------------------------------------------------------
\subsection{Likelihood}\label{sec:model_likelihood}

Conditional on the latent field $\boldsymbol{\eta}$, the observations are independent across $i$.
The conditional probability mass function of $y_i$ given $\eta_i$ (and the inflation parameters) is
\begin{equation}\label{eq:pmf}
  f(y_i \mid \eta_i, \pi_i;\,\kappa) =
  \begin{cases}
    \pi_i + (1-\pi_i)\,h(0;\,\eta_i,\kappa) & \text{if } y_i = 0, \\
    (1-\pi_i)\,h(y_i;\,\eta_i,\kappa) & \text{if } y_i \geq 1,
  \end{cases}
\end{equation}
where
\begin{equation}\label{eq:nb_pmf}
  h(y;\,\eta,\kappa) = \frac{\Gamma(y + \kappa)}{\Gamma(\kappa)\,\Gamma(y+1)}
  \left(\frac{\kappa}{\kappa+e^\eta}\right)^{\!\kappa}
  \left(\frac{e^\eta}{\kappa+e^\eta}\right)^{\!y}
\end{equation}
is the negative binomial probability mass function with mean $\mu = e^\eta$ and shape $\kappa$.

The joint density of $(\by, \boldsymbol{\eta})$ under the restricted ZISAR is
\begin{equation}\label{eq:joint}
  p(\by,\boldsymbol{\eta} \mid \btheta) =
  \left[\prod_{i=1}^{n} f(y_i \mid \eta_i, \pi_i;\,\kappa)\right]
  \cdot p(\boldsymbol{\eta} \mid \bbeta, \rho, \sigma^2),
\end{equation}
where the prior on $\boldsymbol{\eta}$ is the SAR Gaussian:
\begin{equation}\label{eq:sar_prior}
  p(\boldsymbol{\eta} \mid \bbeta,\rho,\sigma^2) =
  (2\pi\sigma^2)^{-n/2}\,|\det(\mathbf{I}_n - \rho\bW)|\,
  \exp\!\left\{-\frac{1}{2\sigma^2}\bigl\|(\mathbf{I}_n - \rho\bW)\boldsymbol{\eta} - \bX\bbeta\bigr\|^2\right\}.
\end{equation}

The marginal log-likelihood, integrating over the latent field, is
\begin{equation}\label{eq:margloglik}
  \ell(\btheta) = \log \int_{\mathbb{R}^n} p(\by,\boldsymbol{\eta} \mid \btheta)\,\mathrm{d}\boldsymbol{\eta}.
\end{equation}
This $n$-dimensional integral has no closed form, but is amenable to Laplace approximation (Section~\ref{sec:estimation}).

% ---------------------------------------------------------------------------
\subsection{Marginal moments}\label{sec:model_moments}

Integrating over $\boldsymbol{\eta}$, the unconditional mean of $y_i$ is
\begin{equation}\label{eq:margmean}
  \E[y_i] = (1 - \pi_i)\,\E[\mu_i] = (1 - \pi_i)\,\E\bigl[e^{\eta_i}\bigr].
\end{equation}
Since $\eta_i$ is the $i$-th component of a multivariate Gaussian (by the reduced form \eqref{eq:reduced_form}), the moment-generating function gives
\begin{equation}\label{eq:E_mu}
  \E\bigl[e^{\eta_i}\bigr] = \exp\!\Bigl(\bar{\eta}_i + \tfrac{1}{2}\,[\Var[\boldsymbol{\eta}]]_{ii}\Bigr),
\end{equation}
where $\bar{\eta}_i = [\mathbf{S}(\rho)\bX\bbeta]_i$ is the conditional mean of $\eta_i$ and $[\Var[\boldsymbol{\eta}]]_{ii} = \sigma^2[\mathbf{S}(\rho)\mathbf{S}(\rho)']_{ii}$.

The marginal variance of $y_i$ is
\begin{equation}\label{eq:margvar}
  \Var[y_i] = (1-\pi_i)\,\E[\mu_i]\bigl(1 + \kappa^{-1}\E[\mu_i]\bigr) + \pi_i(1-\pi_i)\,(\E[\mu_i])^2 + (1-\pi_i)\,\E[\mu_i^2 - (\E[\mu_i])^2],
\end{equation}
which exceeds $\E[y_i]$ whenever $\kappa < \infty$, $\pi_i > 0$, or $\sigma^2 > 0$---confirming that ZISAR accommodates overdispersion through three distinct channels: negative binomial shape ($\kappa$), zero-inflation ($\pi_i$), and spatial random-effects heterogeneity ($\sigma^2$).

% ---------------------------------------------------------------------------
\subsection{Regularity conditions}\label{sec:model_regularity}

\begin{assumption}[Spatial weights]\label{ass:W}
  The sequence of weights matrices $\{\bW_n\}$ is uniformly bounded in absolute row and column sums: $\sup_n \|\bW_n\|_\infty < \bar{c}$ and $\sup_n \|\bW_n\|_1 < \bar{c}$ for some $\bar{c} < \infty$.
  Furthermore, $(\mathbf{I}_n - \rho\bW_n)$ is non-singular for all $\rho \in \mathcal{P}$, where $\mathcal{P}$ is a compact subset of $(-1/r_{\max},\,1/r_{\max})$ and $r_{\max} = \max(|\lambda_1|,|\lambda_n|)$ is the spectral radius of $\bW_n$.
\end{assumption}

\begin{assumption}[Exogenous regressors]\label{ass:X}
  The columns of $\bX$ and $\bZ$ are non-stochastic (or measurable with respect to an exogenous $\sigma$-field) and satisfy $n^{-1}\bX'\bX \to \mathbf{Q}_X$ and $n^{-1}\bZ'\bZ \to \mathbf{Q}_Z$, both positive definite.
\end{assumption}

\begin{assumption}[Parameter space]\label{ass:param}
  The true parameter vector $\btheta_0$ lies in the interior of a compact set $\Theta \subset \mathbb{R}^{\dim(\btheta)}$.
  The shape parameter satisfies $\kappa_0 > 0$, the variance component $\sigma_0^2 > 0$, and $\rho_0 \in \mathcal{P}$.
\end{assumption}

\begin{assumption}[Identification]\label{ass:identification}
  The population marginal log-likelihood $\E[\ell(\btheta)]$ has a unique maximum at $\btheta_0$.
\end{assumption}
